Call a cardinal $\kappa$ $\aleph_0$-inaccessible when $\mu^{\aleph_0}<\kappa$ for every cardinal $\mu<\kappa$. Suppose $\lambda$ is singular and both $\lambda$ and its cofinality $\operatorname{cf}(\lambda)$ have this property. Must
\[\lambda\longrightarrow(\lambda,\aleph_1)^2\]
hold? Thus every red--blue colouring of the unordered pairs of a set of cardinality $\lambda$ would contain either a red homogeneous set of size $\lambda$ or a blue homogeneous set of size $\aleph_1$.
